\section{Related Work}

\subsection{Multimodal Web Agents}
SeeAct \cite{SeeAct} demonstrated that GPT-4V can act as a generalist web agent when supplied with a screenshot and an HTML view, but found grounding to be the dominant bottleneck. AdaptAgent \cite{FewShotLearning:verma2024adaptagentadaptingmultimodalweb} added few-shot demonstrations to extend such agents to unseen websites, while WebExperT \cite{Browsinglikehuman:luo2025browsing} layered fast-and-slow reasoning with experiential memory. Surveys on (M)LLM-based GUI and web agents \cite{tang2025surveymllmbasedguiagents,ning2025surveywebagentsnextgenerationai} consistently flag two open problems: cost and token inefficiency on full-DOM inputs, and weak privacy and safety guarantees inside the action loop.

\subsection{Structural DOM Pruning and Local Grounding}
Real web pages routinely contain $10^4$--$10^5$ DOM tokens \cite{zhang2025prune4web}, well past the practical input window of most LLMs. Prune4Web \cite{zhang2025prune4web} addresses this with a meta-programming filter in which the LLM emits a small keyword/weight scoring script that a deterministic Python scorer applies to the DOM, achieving 25--50$\times$ candidate reduction and 88.28\% Element Accuracy on Mind2Web with a fine-tuned Qwen2.5VL-3B-Instruct grounder. AutoWebGLM \cite{lai2024autowebglm} prunes HTML by trimming non-actionable nodes and limiting tree depth. Agent-E \cite{tarsariya2024agente} compresses the DOM through the accessibility tree and is the only system to monitor cross-step changes via MutationObserver, although it converts changes into prompt text and so does not reduce tokens. WebAgent / HTML-T5 \cite{gur2024webagent} fine-tunes an HTML-specific encoder over long DOM sequences. All four share the property that each step processes a snapshot in isolation.

\subsection{DOM State-Change and Semantic Relevance}
D2Snap \cite{schiepanski2025d2snap} downsamples a DOM to roughly $10^3$ tokens via post-order merging, TextRank-based pruning, and attribute filtering, but operates on static snapshots; Agentic Compilation \cite{agenticcompilation2025} applies per-step DOM sanitisation. The classical algorithmic backdrop for diffing structured trees is Tree Edit Distance \cite{zhangshasha1989ted} with the RTED variant \cite{pawlik2011rted}. None of these are wired into the filter stage of an LLM web agent. On the encoder side, Sentence-BERT \cite{reimers2019sbert} established late fusion of sparse and dense scoring, and MiniLM \cite{wang2020minilm} distils a $\sim$22M-parameter encoder that runs in milliseconds on CPU --- properties that make it a realistic in-loop candidate for relevance scoring on DOM elements. On the grounder side, QLoRA \cite{dettmers2023qlora} makes 4-bit fine-tuning of billion-scale models on a single consumer GPU practical, and the Qwen2.5 \cite{qwen25techreport2024} family supplies permissive 0.5B--0.6B open-weight bases that fit such fine-tuning.

\subsection{Gaps Addressed in This Work}
Three gaps cut across the above. First, sub-1B open-weight grounders are reported only as proof-of-concept variants, without documented training recipes or per-split numbers. Second, all DOM-reduction systems re-process the full snapshot at every step, even when the page has changed only marginally. Third, privacy and HITL safeguards remain bolted on at deployment time rather than designed into the action loop. The remainder of this paper addresses these three gaps directly.
